\notechapter{5}{One-dimensional Bloch, Wannier, and tight-binding warmup}

\textbf{Status: proposed analytical and numerical warmup.} No calculation,
convergence result, fitted model, or plot described in this note has yet been
produced. The purpose is to expose, in a one-dimensional periodic problem, as
many of the representation, alignment, reduction, and convergence issues of
the bulk-silicon program as possible without importing silicon-specific
complexity.

\section*{Purpose and relation to the bulk program}

Consider a particle in a one-dimensional cosine potential. The problem is
simple enough that its Bloch fibers, Fourier couplings, real-space
representation, Wannier transformation, and hopping expansion can all be
written explicitly. It can therefore serve as a controlled analogue of the two
bulk reduction routes:
The direct route is
\begin{equation}
    \text{periodic parent}
    \longrightarrow
    \text{tight binding},
\end{equation}
whereas the operator-mediated route is
\begin{equation}
    \begin{aligned}
        \text{periodic parent}
        &\longrightarrow
        \text{Wannier representation}
        \\
        &\longrightarrow
        \text{tight binding}.
    \end{aligned}
\end{equation}

The warmup should address six questions:
\begin{enumerate}
    \item Do real-space and plane-wave Bloch solvers converge to the same
          retained parent information?
    \item How do basis cutoff, grid spacing, reciprocal sampling, and band
          retention contribute different errors?
    \item How is a localized Wannier representation constructed from the Bloch
          fibers, and which quantities depend on its gauge?
    \item How do real-space hopping coefficients reconstruct the retained Bloch
          operator?
    \item How does a finite hopping range converge to the retained band or
          composite-band operator?
    \item Under which conditions do direct spectral fitting and
          Wannier-mediated hopping truncation give the same reduced model?
\end{enumerate}

The tight-binding model is intended to approximate a retained Bloch operator or
its band dispersion. It should not be described as converging directly to the
scalar cosine potential. Recovering a multiplication potential from reduced
spectral or hopping data is a separate inverse problem.

\section*{Continuum periodic parent}

Let the lattice period be $a$ and define
\begin{equation}
    G=\frac{2\pi}{a}.
\end{equation}
The continuum Hamiltonian on $L^2(\mathbb R)$ is
\begin{equation}
    \hat H
    =
    -\frac{\hbar^2}{2m}\frac{d^2}{dx^2}
    +
    V_0\cos(Gx).
\end{equation}
The natural reciprocal energy scale is
\begin{equation}
    E_G
    =
    \frac{\hbar^2G^2}{2m},
\end{equation}
and the dimensionless potential strength is
\begin{equation}
    \lambda=\frac{V_0}{E_G}.
\end{equation}
The first Brillouin zone may be chosen as
\begin{equation}
    -\frac{G}{2}
    \leq k<\frac{G}{2}.
\end{equation}

By Bloch's theorem, an eigenfunction can be written
\begin{equation}
    \psi_{nk}(x)
    =
    e^{ikx}u_{nk}(x),
    \qquad
    u_{nk}(x+a)=u_{nk}(x).
\end{equation}
The cell-periodic Bloch-fiber operator is
\begin{equation}
    \hat H(k)
    =
    \frac{1}{2m}
    \left(-i\hbar\frac{d}{dx}+\hbar k\right)^2
    +
    V_0\cos(Gx)
\end{equation}
on the periodic cell domain. Equivalently, one may act with the original
differential expression on functions satisfying
\begin{align}
    \psi(a)&=e^{ika}\psi(0),
    \\
    \psi'(a)&=e^{ika}\psi'(0).
\end{align}
The two conventions are unitarily related but must not be mixed within one
matrix construction.

The parent has several exact structural checks:
\begin{align}
    E_n(k+G)&=E_n(k),
    \\
    E_n(-k)&=E_n(k)
\end{align}
for the real inversion-symmetric potential, after consistent band matching.
Changing $V_0$ to $-V_0$ translates the potential by $a/2$ and therefore leaves
the spectrum unchanged. These identities provide useful representation checks
without supplying a convergence result.

\section*{Plane-wave Bloch representation}

Use the Bloch plane waves
\begin{equation}
    \langle x\mid n,k\rangle
    =
    \frac{1}{\sqrt a}
    e^{i(k+nG)x},
    \qquad
    n\in\mathbb Z.
\end{equation}
In this basis the fiber matrix is
\begin{equation}
    H^{\mathrm{PW}}_{nm}(k)
    =
    \frac{\hbar^2(k+nG)^2}{2m}\,\delta_{nm}
    +
    \frac{V_0}{2}
    \left(
        \delta_{n,m+1}+\delta_{n,m-1}
    \right).
\end{equation}
The cosine potential therefore couples only neighboring reciprocal-lattice
indices. This sparsity is an exact consequence of the potential's Fourier
content, not a tight-binding approximation in real space.

For a cutoff $P$, retain
\begin{equation}
    n=-P,\ldots,P,
\end{equation}
which gives a $(2P+1)\times(2P+1)$ Galerkin matrix
$\mathbf H_P^{\mathrm{PW}}(k)$. Increasing $P$ changes the represented parent
space. Low-band convergence should be examined separately from the behavior of
states near the reciprocal cutoff.

The plane-wave matrix makes several analytical limits visible. At $V_0=0$, the
folded free-particle parabolas are exact. For weak nonzero $V_0$, the Fourier
coefficient $V_0/2$ couples crossing free-particle states at a zone boundary.
First-order degenerate perturbation theory predicts a leading gap of
$|V_0|$ at the first coupled crossing. This is expected perturbative behavior
to be checked over an explicitly declared weak-potential regime, not a
numerical result of this note.

\section*{Real-space Bloch-boundary representation}

A second representation acts directly on one spatial cell. Let
\begin{equation}
    x_j=j\Delta x,
    \qquad
    \Delta x=\frac{a}{N},
    \qquad
    j=0,\ldots,N-1.
\end{equation}
The endpoint $x=a$ is identified with $x=0$ through the Bloch phase. With a
centered second-order kinetic representation, define
\begin{equation}
    t_h=\frac{\hbar^2}{2m\Delta x^2}.
\end{equation}
The kinetic matrix has diagonal entries $2t_h$, nearest-neighbor entries
$-t_h$, and boundary entries
\begin{align}
    T^{\mathrm{FD}}_{0,N-1}(k)&=-t_h e^{-ika},
    \\
    T^{\mathrm{FD}}_{N-1,0}(k)&=-t_h e^{ika}.
\end{align}
The potential matrix is
\begin{equation}
    \mathbf V_h
    =
    \operatorname{diag}
    \left(
        V_0\cos(Gx_0),\ldots,
        V_0\cos(Gx_{N-1})
    \right),
\end{equation}
and the represented fiber Hamiltonian is
\begin{equation}
    \mathbf H_h^{\mathrm{FD}}(k)
    =
    \mathbf T_h^{\mathrm{FD}}(k)
    +
    \mathbf V_h.
\end{equation}
This matrix is Hermitian when the two boundary phases are conjugates. Failure
of Hermiticity indicates an inconsistent Bloch closure rather than a physical
band effect.

For a discrete Bloch plane wave with wave number $q=k+nG$, the finite-difference
kinetic eigenvalue is
\begin{equation}
    \varepsilon_h^{\mathrm{kin}}(q)
    =
    \frac{2\hbar^2}{m\Delta x^2}
    \sin^2\left(\frac{q\Delta x}{2}\right).
\end{equation}
For fixed $q$,
\begin{equation}
    \varepsilon_h^{\mathrm{kin}}(q)
    =
    \frac{\hbar^2q^2}{2m}
    \left[
        1-
        \frac{(q\Delta x)^2}{12}
        +O\!\left((q\Delta x)^4\right)
    \right].
\end{equation}
This relation predicts second-order kinetic-dispersion error for fixed smooth
modes. High-frequency grid modes do not satisfy the same small-$q\Delta x$
assumption.

\section*{Direct comparison of the two Bloch representations}

When $N=2P+1$, define the discrete Bloch--Fourier matrix
\begin{equation}
    F_{jn}(k)
    =
    \frac{1}{\sqrt N}
    e^{i(k+nG)x_j},
    \qquad
    n=-P,\ldots,P.
\end{equation}
Its columns are orthonormal in the discrete grid inner product. The
finite-difference Hamiltonian transported into discrete plane-wave coordinates
is
\begin{equation}
    \mathbf A_h^{\mathrm{FD}\rightarrow\mathrm{PW}}(k)
    =
    \mathbf F(k)^\dagger
    \mathbf H_h^{\mathrm{FD}}(k)
    \mathbf F(k).
\end{equation}
Both
\begin{equation}
    \mathbf A_h^{\mathrm{FD}\rightarrow\mathrm{PW}}(k)
    \quad\text{and}\quad
    \mathbf H_P^{\mathrm{PW}}(k)
\end{equation}
then act in the same declared discrete plane-wave coordinates. Their Frobenius
difference is defined:
\begin{equation}
    D_{\mathrm{rep}}(k;h,P)
    =
    \left\lVert
        \mathbf A_h^{\mathrm{FD}\rightarrow\mathrm{PW}}(k)
        -
        \mathbf H_P^{\mathrm{PW}}(k)
    \right\rVert_{\mathrm F}.
\end{equation}

The full-space norm may be dominated by modes near the reciprocal or grid
cutoff. A low-mode comparison should therefore also be defined. Let
$\mathbf P_L$ select reciprocal indices $|n|\leq L<P$. Then
\begin{equation}
    D_{\mathrm{rep}}^{(L)}(k;h,P)
    =
    \left\lVert
        \mathbf P_L
        \left(
            \mathbf A_h^{\mathrm{FD}\rightarrow\mathrm{PW}}(k)
            -
            \mathbf H_P^{\mathrm{PW}}(k)
        \right)
        \mathbf P_L
    \right\rVert_{\mathrm F}
\end{equation}
measures disagreement on a fixed low-reciprocal-mode sector.

The sampled cosine potential also exposes the difference between Galerkin
truncation and cyclic discrete Fourier representation. In an infinite
plane-wave basis, the potential couples $n$ only to $n\pm1$. On a finite
sampled grid, the discrete transform can wrap a coupling across the highest and
lowest represented reciprocal indices. This is an aliasing effect of the
finite cyclic representation. It should not be mistaken for a physical
long-range Fourier coupling.

If two spatial or reciprocal representations have different dimensions, each
must be pulled back to a smaller common Fourier sector or a common retained
band space before a Frobenius norm is formed. Raw subtraction of the unequal
matrices is undefined.

\section*{Mathieu and Floquet formulation}

The cosine problem can also be written as a Mathieu equation. Set
\begin{equation}
    z=\frac{Gx}{2}.
\end{equation}
The stationary Schr\"odinger equation becomes
\begin{equation}
    \frac{d^2\psi}{dz^2}
    +
    \left[
        A-2q\cos(2z)
    \right]\psi
    =0,
\end{equation}
with
\begin{equation}
    A=\frac{4E}{E_G},
    \qquad
    q=\frac{2V_0}{E_G}.
\end{equation}
Bloch periodicity over one lattice period is the Floquet condition
\begin{equation}
    \psi(z+\pi)
    =
    e^{ika}\psi(z).
\end{equation}
Thus the Floquet exponent satisfies
\begin{equation}
    e^{i\pi\nu}=e^{ika},
    \qquad
    \nu=\frac{2k}{G}
    \quad\text{modulo }2.
\end{equation}
Mathieu characteristic values can provide an additional independent reference
for selected bands and crystal momenta. Any such use must state the
characteristic-value convention, branch ordering, Floquet exponent, and energy
conversion. Agreement with one library or table is not assumed by this note.

This formulation is valuable because it separates the differential equation
from either numerical matrix representation. The plane-wave and real-space
solvers can both be checked against the same Floquet problem.

\section*{Band and subspace comparisons}

Let $E_{n,P}^{\mathrm{PW}}(k)$ and $E_{n,h}^{\mathrm{FD}}(k)$ denote matched
band energies. A dimensionless spectral discrepancy is
\begin{equation}
    \delta E_n(k;h,P)
    =
    \frac{
        E_{n,h}^{\mathrm{FD}}(k)
        -
        E_{n,P}^{\mathrm{PW}}(k)
    }{E_G}.
\end{equation}
Band labels must be matched continuously or through subspace information near
close approaches. Sorting eigenvalues independently at every $k$ can hide band
exchange or an inconsistent retained manifold.

If $\mathbf B_{\mathrm{PW}}(k)$ and $\mathbf B_{\mathrm{FD}}(k)$ contain
orthonormal coordinates for corresponding retained subspaces in a common
ambient representation, define
\begin{equation}
    \mathbf S(k)
    =
    \mathbf B_{\mathrm{PW}}(k)^\dagger
    \mathbf B_{\mathrm{FD}}(k).
\end{equation}
The singular values of $\mathbf S(k)$ give principal angles between the
subspaces. A projector discrepancy is
\begin{equation}
    D_P(k)
    =
    \left\lVert
        \mathbf P_{\mathrm{PW}}(k)
        -
        \mathbf P_{\mathrm{FD}}(k)
    \right\rVert_{\mathrm F}.
\end{equation}
Spectral agreement and subspace agreement should be shown separately.

Band-edge curvature supplies a simple observable-level diagnostic. At an
isolated extremum $k_0$, define
\begin{equation}
    \frac{1}{m_n^{\ast}}
    =
    \frac{1}{\hbar^2}
    \left.
        \frac{d^2E_n(k)}{dk^2}
    \right|_{k=k_0}.
\end{equation}
The fit interval, stencil, and reciprocal units must be held fixed when solver,
Wannier, and tight-binding representations are compared.

\section*{Wannier representation and localization}

For an isolated retained band sampled on a uniform mesh of $N_k$ crystal
momenta, let $\mathcal I_{N_k}$ be one complete set of integer residues modulo
$N_k$ and define the $N_k$ inequivalent Born--von Karman lattice vectors
\begin{equation}
    \mathcal R_{N_k}
    =
    \left\{
        R_r=ra
        \;\middle|\;
        r\in\mathcal I_{N_k}
    \right\}.
\end{equation}
Choose phases for the Bloch states and define
\begin{equation}
    \lvert w_R\rangle
    =
    \frac{1}{\sqrt{N_k}}
    \sum_k
    e^{-ikR}
    \lvert\psi_k\rangle,
    \qquad
    R\in\mathcal R_{N_k}.
\end{equation}
A smooth periodic phase change modifies the Wannier profile but leaves its
center modulo $a$ fixed by the isolated-band Berry phase. A gauge with nonzero
reciprocal winding can shift the representative center by a lattice vector.
The isolated-band projector remains invariant. In one dimension, this freedom
provides a direct setting in which to study gauge continuity, lattice-equivalent
center shifts, spread, Berry phase, and the relation between reciprocal
regularity and real-space decay~\cite{kohn1959,marzari2012}. In the dense-mesh
limit, the Born--von Karman representatives extend to the ordinary infinite
lattice.

\subsection*{Direct isolated-band parallel transport}

Let $\lvert u_j\rangle$ be normalized cell-periodic states at neighboring mesh
points $k_j$ and $k_{j+1}$. Define
\begin{equation}
    m_j
    =
    \langle u_j\mid u_{j+1}\rangle.
\end{equation}
When $m_j\neq0$, multiply $\lvert u_{j+1}\rangle$ by a phase that makes $m_j$
real and nonnegative. Repeating this operation along the reciprocal mesh gives
a parallel-transport gauge. The accumulated phase on closing the Brillouin
zone is the discrete holonomy. It must be distributed consistently across the
mesh to obtain a periodic gauge.

This construction is sufficiently transparent that every phase change can be
inspected. The resulting Wannier center is related to the accumulated Berry or
Zak phase modulo a lattice vector. Translating the Wannier function by one cell
changes the representative center but not the physical polarization class or
the retained band dispersion. The origin and phase convention must therefore
be retained with every center comparison.

A vanishing or very small $|m_j|$ signals that neighboring states have not been
matched reliably, that the mesh is too coarse, or that an isolated-band
description is failing near a crossing. Phase smoothing cannot repair a loss of
subspace overlap.

\subsection*{Composite-band parallel transport}

For an $M$-band retained space, define the neighboring overlap matrix
\begin{equation}
    \mathbf M^{(j)}_{mn}
    =
    \langle u_{m,k_j}\mid u_{n,k_{j+1}}\rangle.
\end{equation}
Its polar or singular-value decomposition supplies the unitary rotation that
maximizes neighboring-frame overlap. Sequential application produces a
parallel-transported composite frame. The Wilson product around the reciprocal
mesh records the closure holonomy, whose eigenphases determine the centers of
hybridized one-dimensional Wannier channels modulo lattice translations.

Small singular values of $\mathbf M^{(j)}$ diagnose incompatible neighboring
subspaces. They cannot be corrected by a unitary gauge rotation. This provides
a controlled analogue of the principal-angle and alignment checks required for
the silicon retained manifold.

\subsection*{Wannier90 as an independent localization route}

Wannier90 can localize Bloch states that originate from this model; the parent
states need not come from a density-functional calculation. It consumes the
Bloch eigenvalues and the neighboring-state overlaps
\begin{equation}
    M_{mn}^{(\mathbf k,\mathbf b)}
    =
    \langle
        u_{m\mathbf k}
        \mid
        u_{n,\mathbf k+\mathbf b}
    \rangle,
\end{equation}
together with optional trial-orbital projections
\begin{equation}
    A_{mn}(\mathbf k)
    =
    \langle
        \psi_{m\mathbf k}
        \mid
        g_n
    \rangle.
\end{equation}
For the model experiment, the retained interface consists of the reciprocal
mesh and neighbor list, eigenvalues, overlap matrices, optional projections,
lattice convention, band count, and Wannier count. In the standard Wannier90
file interface these data are represented by the \texttt{.eig},
\texttt{.mmn}, optional \texttt{.amn}, and \texttt{.win} surfaces. Wannier90
then performs gauge selection, optional disentanglement, localization, and
interpolation~\cite{souza2001,pizzi2020}.

The difficult interface quantity is the overlap matrix, not the scalar band
energy. It must use cell-periodic states, the correct represented inner
product, and consistent reciprocal-boundary sewing. For real-space states the
inner product includes the declared quadrature. For plane-wave states it uses
the coefficient inner product with the correct reciprocal relabelling at a
Brillouin-zone boundary. Arbitrary phases from independent eigensolutions are
acceptable only when the overlap matrices are constructed covariantly from
those same states.

A one-dimensional model may be represented for Wannier90 with one reciprocal
point in each inactive direction and a declared auxiliary transverse cell and
orbital factor. The transverse lattice and spread contribution are then
representation conventions, not physical dimensions of the model. Active-axis
centers and spreads should be reported separately from any auxiliary transverse
contribution.

\subsection*{Comparison of the two localization routes}

The direct parallel-transport construction and Wannier90 should consume the
same retained Bloch states and reciprocal mesh. Their outputs should be
compared only after lattice translations, orbital order, phases, and any
composite-space unitary freedom have been aligned. Proposed diagnostics are:
\begin{enumerate}
    \item retained-projector agreement;
    \item Wannier centers modulo $a$;
    \item active-axis spreads;
    \item real-space profiles after translation and phase alignment;
    \item reconstructed band interpolation;
    \item aligned hopping coefficients; and
    \item sensitivity to reciprocal-mesh density and trial projections.
\end{enumerate}
Agreement would verify two representations of the declared localization
problem. Disagreement should first be attributed among overlap construction,
reciprocal sewing, gauge closure, projection choice, disentanglement, and
Fourier convention before it is interpreted as a localization failure.

\subsection*{Hamiltonian in Wannier coordinates}

For one isolated band, the represented Hamiltonian hopping coefficients are
\begin{equation}
    t(R)
    =
    \langle w_0\rvert\hat H\lvert w_R\rangle
    =
    \frac{1}{N_k}
    \sum_k
    e^{-ikR}E(k),
    \qquad
    R\in\mathcal R_{N_k}.
\end{equation}
The inverse discrete transform is
\begin{equation}
    E(k)
    =
    \sum_{R\in\mathcal R_{N_k}}
    e^{ikR}t(R).
\end{equation}
For an isolated band, these scalar hoppings depend only on the represented band
dispersion, even though the Wannier functions themselves depend on phase.

For a composite group of bands, the retained Hamiltonian is matrix valued:
\begin{equation}
    \mathbf H_{\mathrm W}(R)
    =
    \frac{1}{N_k}
    \sum_k
    e^{-ikR}
    \mathbf H_{\mathrm W}(k),
    \qquad
    R\in\mathcal R_{N_k}.
\end{equation}
A momentum-dependent unitary gauge changes the individual matrix blocks and
Wannier functions. Operator comparisons then require a common gauge or an
explicit alignment map, as in the bulk-silicon problem.

\section*{Finite-range tight-binding reduction}

For an isolated band, choose a symmetric retained set
\begin{equation}
    \mathcal R_c
    \subseteq
    \mathcal R_{N_k}
\end{equation}
of distinct Born--von Karman representatives whose physical distances do not
exceed the declared cutoff. The finite-range dispersion is
\begin{equation}
    E_{\mathrm{TB}}^{(\mathcal R_c)}(k)
    =
    \sum_{R\in\mathcal R_c}
    e^{ikR}t(R).
\end{equation}
Hermiticity requires that $\mathcal R_c$ be closed under inversion modulo the
Born--von Karman superlattice and that
\begin{equation}
    t(-R)=t(R)^{\ast}
\end{equation}
under the selected representative convention. For the real
inversion-symmetric problem, a centered representative set and compatible
origin give the usual cosine expansion when the hoppings are real.

The omitted-hopping residual is
\begin{equation}
    \Delta E^{(\mathcal R_c)}(k)
    =
    E(k)-E_{\mathrm{TB}}^{(\mathcal R_c)}(k)
    =
    \sum_{R\in\mathcal R_{N_k}\setminus\mathcal R_c}
    e^{ikR}t(R).
\end{equation}
On a complete uniform discrete transform, Parseval's identity gives
\begin{equation}
    \sum_k
    \left|
        \Delta E^{(\mathcal R_c)}(k)
    \right|^2
    =
    N_k
    \sum_{R\in\mathcal R_{N_k}\setminus\mathcal R_c}
    |t(R)|^2.
\end{equation}
Thus the root-sum-square spectral error and the omitted hopping norm are two
coordinate descriptions of the same isolated-band truncation residual under
these specific finite-transform conventions.

For a composite retained manifold, define
\begin{equation}
    \mathbf H_{\mathrm{TB}}^{(\mathcal R_c)}(k)
    =
    \sum_{R\in\mathcal R_c}
    e^{ikR}\mathbf H_{\mathrm W}(R).
\end{equation}
The corresponding operator residual is
\begin{equation}
    \Delta\mathbf H^{(\mathcal R_c)}(k)
    =
    \mathbf H_{\mathrm W}(k)
    -
    \mathbf H_{\mathrm{TB}}^{(\mathcal R_c)}(k).
\end{equation}
This is the one-dimensional analogue of truncating Wannier Hamiltonian blocks
by neighbor range in bulk silicon.

\section*{Direct and Wannier-mediated reduction routes}

The direct route fits a finite-range scalar dispersion or matrix-valued model
to selected parent band data. For one isolated band, a weighted direct
objective may be written
\begin{equation}
    \mathcal L_{\mathrm{dir}}
    \left(\{t(R)\}_{R\in\mathcal R_c}\right)
    =
    \sum_{k\in\mathcal T}
    w_k
    \left|
        E_{\mathrm{parent}}(k)
        -
        \sum_{R\in\mathcal R_c}e^{ikR}t(R)
    \right|^2.
\end{equation}
The mediated route first computes the complete set of $N_k$ represented
hopping coefficients through the inverse discrete Fourier transform and then
retains $\mathcal R_c$.

For a complete uniform $k$ grid, equal weights, a consistent discrete Fourier
convention, and an unconstrained Hermitian hopping class whose retained lattice
vectors are distinct modulo $N_k$, the two procedures are expected to agree:
both are the orthogonal projection of the represented dispersion onto the same
finite Fourier subspace. This is an algebraic property
to verify, not an empirical coincidence.

The routes need not agree when:
\begin{itemize}
    \item the direct fit uses nonuniform reciprocal weights;
    \item training and withheld $k$ points differ;
    \item symmetry or parameter-sharing constraints restrict the direct class;
    \item the Wannier route uses a composite-band gauge or alignment;
    \item the reciprocal and real-space transforms use incompatible grids or
          normalization;
    \item optimization does not find the declared minimum; or
    \item additional observables enter one objective but not the other.
\end{itemize}

Let $\mathbf H_{\mathrm{TB,dir}}(k)$ and
$\mathbf H_{\mathrm{TB,med}}(k)$ be final aligned operators. The route defect is
\begin{equation}
    D_{\mathrm{path}}^2
    =
    \sum_{k\in\mathcal K_{\mathrm{cmp}}}
    w_k
    \left\lVert
        \mathbf H_{\mathrm{TB,dir}}(k)
        -
        \mathbf H_{\mathrm{TB,med}}(k)
    \right\rVert_{\mathrm F}^2.
\end{equation}
The scalar isolated-band case supplies a setting in which a zero route defect
can be expected under precisely stated ideal conditions. Controlled departures
from those conditions can then show how path dependence enters before the more
complicated silicon problem is attempted.

\section*{What can be learned from varying the cosine strength}

The dimensionless amplitude $\lambda=V_0/E_G$ connects qualitatively different
regimes within one parent family:
\begin{description}
    \item[Free-particle limit.] At $\lambda=0$, the plane-wave basis is exact,
      the folded parabolas cross, and localized single-band constructions at
      crossings require care.
    \item[Weak-potential regime.] Small $|\lambda|$ opens gaps at coupled
      crossings and tests nearly-free-particle perturbative expectations.
    \item[Intermediate regime.] Bands become isolated over wider energy
      ranges, allowing gauge, Wannier, and hopping-decay studies without an
      extreme asymptotic approximation.
    \item[Deep-potential regime.] Low bands narrow and localized-orbital
      descriptions are expected to become more effective, while resolving the
      sharply varying cell wavefunctions may require larger plane-wave cutoffs
      or finer real-space grids.
\end{description}

The sign change $V_0\mapsto -V_0$ provides a translation-equivalence check. A
representation that respects the half-cell translation should reproduce the
same band energies while shifting appropriate real-space centers. This separates
spectral invariance from coordinate-dependent Wannier data.

The amplitude scan can also test how hopping decay, required TB range, band
gaps, bandwidths, effective masses, and Wannier spreads change together. These
relationships are proposed observations; they must not be asserted before the
calculations are performed.

\section*{Proposed staged experiment}

\subsection*{Stage A: parent representation verification}

Choose a frozen set of dimensionless amplitudes $\lambda$ and crystal momenta.
For each amplitude:
\begin{enumerate}
    \item refine the plane-wave cutoff $P$ at fixed $k$ points;
    \item refine the real-space grid spacing $\Delta x$ at the same points;
    \item compare low-mode operator blocks after Fourier transport;
    \item compare retained eigenvalues, projectors, and symmetry identities;
    \item evaluate selected Mathieu characteristic values as an independent
          Floquet reference when their conventions are controlled; and
    \item freeze a parent representation only after the declared quantities
          are stable under both refinement routes.
\end{enumerate}

The plane-wave and real-space errors must be estimated independently. Agreement
between two unconverged representations is not sufficient.

\subsection*{Stage B: reciprocal-space sampling}

Using the accepted parent representation, refine the $k$ mesh independently of
the plane-wave or spatial representation. Examine:
\begin{enumerate}
    \item band interpolation between sampled points;
    \item zone-center and zone-boundary energies;
    \item the first band gap and selected higher gaps;
    \item bandwidths and band-edge curvatures; and
    \item stability of Fourier-derived hopping coefficients.
\end{enumerate}
A dense $k$ mesh cannot repair an unconverged Bloch-fiber representation, and a
converged fiber solver does not guarantee converged real-space hopping decay on
an insufficient reciprocal mesh.

\subsection*{Stage C: direct isolated-band localization}

Select an isolated band over the full Brillouin zone. Construct its gauge by
discrete parallel transport, record the neighboring overlaps, close the gauge
periodically, and retain the accumulated Berry phase. Report the center modulo
$a$, active-axis spread, and real-space profile. Verify that the
Fourier-transformed hoppings reconstruct the represented band on the original
$k$ mesh.

Repeat the construction under controlled phase changes. The band energies and
isolated-band hoppings should remain invariant under compatible phase changes,
while the Wannier function may translate or change its coordinate
representation. This separates gauge-dependent localization data from the
retained scalar operator.

\subsection*{Stage D: independent Wannier90 localization}

Construct the neighboring overlap matrices and optional trial-orbital
projections from the same retained Bloch states. Retain the reciprocal sewing,
inner-product convention, lattice embedding, and active versus inactive
spatial directions. Use Wannier90 to obtain an independently localized and
interpolated representation.

Align the direct and Wannier90 results and compare their projectors, centers
modulo $a$, active-axis spreads, profiles, interpolation, and hopping
coefficients. Repeat over reciprocal-mesh refinements and controlled projection
choices. This stage tests the localization interface and gauge construction; it
does not treat agreement between two outputs as validation of the cosine model
against experiment.

\subsection*{Stage E: composite-band construction}

Retain two or more bands and introduce a momentum-dependent unitary frame.
Construct both a direct polar-transport frame and a Wannier90 frame from the
same overlap data. This stage should examine:
\begin{enumerate}
    \item projector invariance under gauge rotation;
    \item singular values of neighboring overlap matrices;
    \item Wilson-loop closure and center phases;
    \item changes in individual Wannier functions and matrix blocks;
    \item covariance of the reciprocal Hamiltonian;
    \item alignment recovery after a controlled gauge transformation; and
    \item the effect of gauge smoothness on real-space block decay.
\end{enumerate}
The composite stage is the closer analogue of the multiorbital
$sp^3s^{\ast}$ silicon problem.

\subsection*{Stage F: hopping-range hierarchy}

For each selected isolated or composite retained space, construct a frozen
hierarchy
\begin{equation}
    r_c=0,1,2,\ldots,r_{\max}.
\end{equation}
At each level report:
\begin{enumerate}
    \item retained and omitted hopping norms;
    \item reciprocal-space Frobenius residuals;
    \item band-energy errors over training and withheld points;
    \item gap, bandwidth, and effective-mass errors; and
    \item sensitivity to the Wannier gauge in the composite case.
\end{enumerate}
The lowest range meeting every prespecified criterion is the selected reduced
class for the warmup. If none passes, the untruncated represented operator
remains the reference.

\subsection*{Stage G: direct versus mediated TB routes}

Fit the same declared finite-range class directly to frozen parent band data,
then compare it with the hopping-truncated Wannier result. Begin with the
isolated-band, uniform-grid, equal-weight case where route agreement is
expected algebraically. Introduce one controlled difference at a time:
\begin{enumerate}
    \item nonuniform reciprocal weights;
    \item withheld reciprocal regions;
    \item shared or symmetry-constrained hopping parameters;
    \item a composite retained space;
    \item independent gauge choices requiring alignment; and
    \item simultaneous spectral and operator objectives.
\end{enumerate}
This sequence turns path dependence into an observable controlled phenomenon
rather than an unexplained disagreement between two fitted models.

\section*{Error accounting}

The warmup should retain at least the following components:
\begin{enumerate}
    \item plane-wave cutoff error;
    \item real-space grid and stencil error;
    \item Bloch-boundary implementation error, assessed through exact
          structural checks;
    \item reciprocal $k$-mesh and Fourier-transform error;
    \item retained-band or retained-subspace error;
    \item neighboring-overlap, reciprocal-sewing, and Gram-conditioning error;
    \item direct and Wannier90 localization-route discrepancy;
    \item gauge and alignment sensitivity;
    \item hopping-range truncation error;
    \item direct-fit optimization and identifiability error; and
    \item observable-extraction error for gaps, bandwidths, and curvatures.
\end{enumerate}
These quantities should not be combined into one total error without a declared
relationship among their norms, units, and dependencies. The cosine parent is
an illustrative mathematical model, so agreement within this experiment is
numerical verification of the declared construction rather than scientific
validation of silicon.

\section*{Plot placeholders}

\begin{enumerate}
    \item \textbf{Placeholder---potential and unit cell.} Plot
          $V_0\cos(Gx)/E_G$ over several cells for the selected amplitudes,
          marking the cell origin and translation relating $V_0$ to $-V_0$.
    \item \textbf{Placeholder---Bloch bands from two parent representations.}
          Overlay plane-wave and real-space bands over the first Brillouin zone
          for several refinement levels.
    \item \textbf{Placeholder---Mathieu comparison.} Plot selected
          plane-wave, real-space, and Mathieu/Floquet energies after applying a
          common branch and unit convention.
    \item \textbf{Placeholder---plane-wave convergence.} Plot low-band energy
          and projector discrepancies against reciprocal cutoff $P$.
    \item \textbf{Placeholder---real-space convergence.} Plot the same
          quantities against $\Delta x/a$ at fixed potential strength.
    \item \textbf{Placeholder---representation Frobenius discrepancy.} Plot
          $D_{\mathrm{rep}}^{(L)}(k;h,P)/E_G$ against $k$, $P$, and
          $\Delta x/a$ for fixed low-mode sectors.
    \item \textbf{Placeholder---matrix structure.} Display heat maps of the
          plane-wave Galerkin matrix and the Fourier-transported real-space
          matrix, highlighting kinetic-dispersion differences and aliasing at
          the cutoff.
    \item \textbf{Placeholder---symmetry checks.} Plot residuals for
          $E_n(k)-E_n(-k)$, reciprocal periodicity, and the
          $V_0\leftrightarrow -V_0$ translation equivalence.
    \item \textbf{Placeholder---gap opening.} Plot the first zone-boundary gap
          against $|V_0|/E_G$, including the declared weak-potential comparison
          with the leading perturbative expectation.
    \item \textbf{Placeholder---Bloch states.} Plot the density and phase of
          selected Bloch states at the zone center, zone boundary, and an
          interior $k$ point.
    \item \textbf{Placeholder---Wannier functions.} Plot isolated- and
          composite-band Wannier functions over several cells for controlled
          gauge choices.
    \item \textbf{Placeholder---localization-route comparison.} Overlay the
          aligned direct parallel-transport and Wannier90 Wannier functions and
          compare their interpolated bands and hopping blocks.
    \item \textbf{Placeholder---neighbor-overlap diagnostics.} Plot isolated
          overlap magnitudes, composite overlap singular values, and closure
          phases over the reciprocal mesh.
    \item \textbf{Placeholder---Wannier centers and spreads.} Plot direct and
          Wannier90 centers and active-axis spreads against potential strength,
          reciprocal-mesh density, and gauge or projection convention.
    \item \textbf{Placeholder---hopping decay.} Plot $|t(R)|/E_G$ or
          $\lVert\mathbf H_{\mathrm W}(R)\rVert_{\mathrm F}/E_G$ against
          $|R|/a$ for several bands and potential strengths.
    \item \textbf{Placeholder---TB range convergence.} Plot spectral and
          operator discrepancies separately against hopping range $r_c$.
    \item \textbf{Placeholder---Parseval check.} Compare the reciprocal-space
          residual norm with the omitted-hopping norm under a complete uniform
          transform.
    \item \textbf{Placeholder---direct and mediated routes.} Overlay direct-fit
          and Wannier-mediated TB bands and plot their route defect against
          hopping range.
    \item \textbf{Placeholder---controlled path dependence.} Plot
          $D_{\mathrm{path}}$ as reciprocal weights, fitting domains, gauge
          choices, and model constraints are varied one at a time.
    \item \textbf{Placeholder---effective mass and bandwidth.} Plot errors in
          selected band-edge curvatures and bandwidths against parent
          representation cutoff, $k$-mesh density, and hopping range.
    \item \textbf{Placeholder---two-parameter convergence maps.} Display
          selected errors over $(P,\Delta x/a)$, $(N_k,r_c)$, and
          $(|V_0|/E_G,r_c)$ while keeping the remaining controls fixed.
\end{enumerate}
Each completed plot should identify the potential strength, solver
representation, cutoff or grid, reciprocal mesh, retained bands, gauge,
alignment, hopping range, normalization, units, and evidentiary status. A plot
placeholder is not a calculated result.

\section*{Interpretation boundary}

This warmup is designed to make the full reduction chain visible in a model
whose Fourier structure can be inspected directly. The parent comparison tests
whether independent Bloch representations agree after alignment. The direct and Wannier90 localization routes separate overlap construction,
subspace selection, gauge choice, and localization from operator reconstruction.
The hopping hierarchy turns real-space truncation into a controlled model
reduction. The direct and mediated routes expose when spectral fitting and
operator truncation commute and when they do not.

Successful completion would provide analytical and numerical verification for
the declared one-dimensional construction. It would not establish that the
same convergence rates, gauges, hopping ranges, or model classes apply to
bulk silicon. Its value is methodological: every ambiguity encountered here can
be resolved in a setting where the state spaces, Fourier coefficients, and
limiting operations are explicit before the corresponding silicon comparison
is attempted.
